% Taxonomie des réponses. Les catégories sont des états, pas des séries :
% la couleur suit donc une échelle de statut, du vert (résultat visé) au
% rouge (échec), et chaque barre porte son libellé — l'identité ne repose
% jamais sur la seule couleur.
%
% Deux réglages non évidents :
%   - les libellés d'axe sont écrits en clair : pgfplots ne développe pas une
%     macro passée à « yticklabels », quelle que soit la manière de la lui
%     donner. Leur ORDRE doit rester celui de TAXONOMY_ORDER dans
%     rapport/build_data.py, qui fixe l'ordre des barres ;
%   - « bar shift=0 » : en mode xbar, chaque \addplot reçoit sinon un décalage
%     propre et les barres se retrouvent échelonnées au lieu d'être alignées
%     sur leur ordonnée.
% Taxonomie des réponses générées, du meilleur au pire résultat
% Généré par rapport/build_data.py — NE PAS ÉDITER.
\newcommand{\TaxonomieLabels}{{Sourcée, article attendu},{Sourcée, autre article},{Abstention correcte},{Citation non vérifiée},{Sans citation},{Refus à tort}}
\newcommand{\TaxonomieBarA}{(41,0)}
\newcommand{\TaxonomieBarB}{(2,1)}
\newcommand{\TaxonomieBarC}{(9,2)}
\newcommand{\TaxonomieBarD}{(6,3)}
\newcommand{\TaxonomieBarE}{(2,4)}
\newcommand{\TaxonomieBarF}{(2,5)}
\newcommand{\TaxonomieMax}{41}
\newcommand{\TaxonomieTotal}{62}

\begin{tikzpicture}
\begin{axis}[
  barresH,
  height=6cm,
  % La largeur de pgfplots ne compte que l'aire tracée : les libellés de
  % catégorie s'y ajoutent. On la réduit pour que l'ensemble tienne dans la
  % justification.
  width=0.74\linewidth,
  xmin=0, xmax=\TaxonomieMax*1.16,
  ytick={0,...,5},
  yticklabels={{Sourcée, article attendu},{Sourcée, autre article},
               {Abstention correcte},{Citation non vérifiée},
               {Sans citation},{Refus à tort}},
  yticklabel style={font=\sffamily\scriptsize, text=ink, align=right},
  xlabel={nombre de questions \textcolor{slate}{(sur \TaxonomieTotal)}},
  nodes near coords,
  bar width=11pt,
  every axis plot/.append style={bar shift=0},
]
\addplot[draw=leafgreen, fill=leafgreen, fill opacity=0.88] coordinates {\TaxonomieBarA};
\addplot[draw=leafgreen, fill=leafgreen, fill opacity=0.45] coordinates {\TaxonomieBarB};
\addplot[draw=inptblue,  fill=inptblue,  fill opacity=0.85] coordinates {\TaxonomieBarC};
\addplot[draw=amberdark, fill=amberdark, fill opacity=0.85] coordinates {\TaxonomieBarD};
\addplot[draw=amberdark, fill=amberdark, fill opacity=0.5]  coordinates {\TaxonomieBarE};
\addplot[draw=anrtred,   fill=anrtred,   fill opacity=0.85] coordinates {\TaxonomieBarF};
\end{axis}

% La catégorie absente est le vrai résultat : aucune réponse produite hors
% corpus. Elle est posée sous la figure, hors de l'axe, pour ne pas se
% confondre avec une barre de valeur nulle.
\node[font=\sffamily\scriptsize, text=leafgreen!40!black, anchor=north west,
      align=left, inner sep=0pt]
  at ($(current bounding box.south west)+(0,-3mm)$)
  {\textbf{Répond alors qu'il devrait s'abstenir : 0} \textcolor{slate}{— catégorie vide}};
\end{tikzpicture}
